%\documentclass[conference, 10pt]{llncs}
\documentclass[conference, 10pt]{IEEEtran}
\usepackage{graphicx} % Required for inserting images

%\title{Hacia una Inteligencia Artificial Eficiente: Uso de FPGAs para la Sostenibilidad Energética}
%\title{Towards an efficient Artificial Inteligence: Use of FPGAs for energetic sustainabiity}
\title{Towards a efficient Artificial inteligence: Use of FPGAs for energetic sustainavility}
\author{Unai Fernandez Urroz}
\date{December 2025}

\begin{document}

\maketitle

\begin{abstract}

La creciente complejidad de las técnicas de Inteligencia Artificial y el uso de conjuntos de 
datos cada vez más grandes, ha llevado a un aumento significativo en cuanto al consumo 
energético en implementaciones en hardware convencional como CPU y GPU.  

%%Este desafío resalta la necesidad de enfoques más sostenibles, como la Green AI, que busca optimizar el uso de recursos y reducir el impacto ambiental.

El uso de FPGAs (Field Programable Gate Arrays) en este contexto surge como una alternativa prometedora, debido a las
ventajas que ofrecen. La capacidad de diseñar aceleradores personalizados permite mejorar tanto la latencia como el
rendimiento del sistema, manteniendo una eficiencia energética excelente en todo momento. Esta eficiencia energética
obtenida, facilita la implementación de algoritmos de aprendizaje automático en sistemas embebidos con recursos
limitados.

El principal objetivo de este trabajo es implementar la inferencia de conjuntos de arboles de 
decisión utilizando el algoritmo XGBoost en FPGA buscando una mejor eficiencia energética, 
manteniendo un buen rendimiento, comparada con otras implementaciones en FPGA y hardware convencional.

%El principal objetivo de la tesis es implementar el entrenamiento y la inferencia de algoritmos clásicos de aprendizaje automático en FPGA buscando una mejor eficiencia energética manteniendo un buen rendimiento. Las arquitecturas desarrolladas se compararán con otras propuestas en FPGA y con los resultados obtenidos en hardware convencional. Finalmente se buscará integrar los diseños de hardware implementados en un sistema embebido para evaluar su efectividad en un entorno real.

\end{abstract}
\section{Introduction}

%todo%
\begin{itemize}
	\item Hacer una pequeña introducción sobre IA y su impacto en los ultimos años. Hacer mas incapie en arboles de decision.
	\item Mecionar por que es viable usar FPGA para acelerar inferencia y entrenamiento de algoritmos de ML
	\item Estado del arte sobre el tema. Comentar varios trabajos sobre arboles de decisiones y algun otro acelerador para otro algoritmo que pueda ser interesante.
	\item Las partes que tiene el documento.
\end{itemize}

The area of artificial Intelligence (AI) focuses on the development of systems capable of \textbf{imitating} the human brain ability of decision taking and resolution of problems.In recent years the popularity of AI ahs risen considerably ampliando su presencia a muchos ambitos como la industria, a medicina o incluso la educacion.      


\section{Background}
\section{Design and implemtentation of the hardware accelerator}
\subsection{Decision tree accelerator}
\subsection{Main structure}
\section{Training procedure}
\subsection{Python app}
\section{Evaluation of the designed accelerator}
\section{Conclusions}

La Inteligencia Artificial (IA) es un área de la informática que se enfoca en el diseño de sistemas que pueden imitar
la capacidad de toma de decisiones y resolución de problemas del cerebro humano. Los últimos años la popularidad
de la IA ha aumentado considerablemente y está presente en muchos ámbitos, como la industria, la medicina o
incluso la educación. Una de las ramas dentro de la IA es el aprendizaje automático (Machine Learning, ML), cuyo
propósito es que los sistemas informáticos aprendan a realizar una tarea especifica basándose en una serie de datos
analizados. Desde hace varias décadas se han desarrollado diversos algoritmos, con varios niveles de complejidad e
idoneidad para diferentes aplicaciones dentro de esta área. Algunos de los algoritmos más conocidos son los árboles
de decisión (Decision Tree, DT), Random Forest (RF), Support Vector Machine (SVM) y las redes neuronales
convolucionales (CNN).

La creciente complejidad de estas técnicas junto con el uso de conjuntos de datos más cada vez grandes aumenta el
consumo energético al implementarlas en harware convencional como CPUs o GPUs. A lo largo de los años, el
hardware ha experimentado una notable evolución, lo que ha permitido mejorar diversos aspectos de rendimiento
y eficiencia. Sin embargo, aún existen oportunidades para seguir optimizando estas áreas. En este contexto, las
FPGAs (Field Programmable Gate Array) emergen como una alternativa que reduce tanto el consumo energético
como la latencia del sistema. Además, permiten la creación de sistemas digitales totalmente personalizados, lo que
permite optimizar el uso de los recursos disponibles. Esta característica y el bajo consumo energético abre la puerta
a la integración de algoritmos de ML en sistemas embebidos, lo que posibilita el desarrollo de una amplia variedad
de nuevas aplicaciones.

El estado del arte actual presenta una gran variedad de propuestas utilizando diferentes modelos de ML, y con
aplicaciones específicas en diversos campos. Olmos et al. \cite{olmosoptimized} proponen la implementación de una arquitectura en
una FPGA, para optimizar el rendimiento del algoritmo XGBoost. Por otro lado, Owaida y colaboradores \cite{owaida2017scalable}
presentan una arquitectura para procesar varios árboles de decisión, en la que se tiene en cuenta el mapeo en
memoria de los datos y el acceso a ellos.

Shabarinath et al. \cite{10969767} por su parte, presentan una arquitectura para entrenar un modelo SVM usando el algoritmo
SMO (Secuential Minimal Optimization) \cite{10.5555/299094.299105}. Además, la arquitectura también incluye una función para hacer
inferencia. En otros trabajos como el de Yacoub y su equipo \cite{YACOUB2024100115} se implementan dos técnicas de generalización (OvA
y OvO) para el algoritmo SVM, para lidiar con los problemas de multi-class. La representación de los datos juega un
papel importante para aprovechar los recursos de memoria, Zhou et al. \cite{7490869} proponen una solución para acelerar el
reconocimiento de imágenes con redes neuronales profundas, en la que se usa fixed-point con 11 bits de precisión
para representar los números reales.

La implementación de algoritmos de IA en FPGA puede ser compleja, y muchas de las veces se opta por realizar el
entrenamiento con herramientas ya conocidas, e implementar los aceleradores en FPGA para la inferencia. Otro
aspecto importante para considerar es la gestión de la memoria. Almacenar los datos necesarios para la ejecución
del modelo y aprovechar al máximo el espacio de almacenamiento puede resultar complejo. Además, la
representación de los datos tiene un impacto significativo en la eficiencia del uso de la memoria, lo que puede afectar el rendimiento general del sistema.

\bibliographystyle{ieeetr}
\bibliography{references.bib}

\end{document}
